
\chapter{Санкционная проверка}\label{ch:sanctions}

\highlight{У~санкционной проверки (\textit{sanctions screening}) собственные
  перечни, алгоритмы сравнения, пороги ложных срабатываний и~маршрут эскалации;
к~ПОД/ФТ она~не сводится}. ОПДС работает с~национальными перечнями
Росфинмониторинга и~решениями Межведомственной комиссии (МВК)
при~Росфинмониторинге; международные перечни (Управление по~контролю
  над~иностранными активами Министерства финансов США, \enquote{Office of Foreign
Assets Control}, OFAC; Европейский союз, EU; Организация Объединённых Наций, UN)
применяются по~собственной политике ОПДС и~требованиям банков-корреспондентов.
Если клиент, бенефициар или~контрагент попадает в~применимый перечень, операцию
останавливают до~исполнения.

Российский банк не~может ограничиться национальными перечнями, даже если
сам в~SDN-листе США не~значится. Любой трансграничный расчёт проходит через
банк-корреспондент в~валюте перевода; ОПДС проверяет исходящие платежи,
а~банк-корреспондент повторно сверяет входящие сообщения с~OFAC SDN, EU
Consolidated и~UK Sanctions List и~обязан заблокировать средства
при~совпадении\footnote{Именно входящая проверка на~стороне корреспондента
  закрывает канал отправителю: миновать фильтр получателя физически невозможно.
Технические попытки обойти его описаны в~разделе~\ref{sec:sanctions-wire-stripping}.}.
Пропуск рискует не~только остановкой одной операции, но~и~разрывом корреспондентской
линии целиком. Конкретный механизм давления на~корреспондентов разобран
в~разделе~\ref{sec:sanctions-international}.

\section{Национальные перечни}\label{sec:sanctions-rf-lists}

Для российского ОПДС~--- банка или~НКО (см.~гл.~\ref{ch:regulation})~--- первичны
три национальных источника. Главный~--- перечень организаций и~физических лиц,
в~отношении которых имеются сведения об~их причастности к~экстремистской
деятельности или~терроризму; ведёт его Росфинмониторинг, а~меры по~замораживанию
средств применяются по~подпункту~6 пункта~1 статьи~7 115-ФЗ (см.~гл.~\ref{ch:aml};
  операции с~лицами из~перечня дополнительно подлежат обязательному контролю
по~статье~6). Параллельно ведётся перечень лиц, в~отношении которых имеются
сведения об~их причастности к~распространению оружия массового уничтожения.
Третий источник~--- решения МВК о~замораживании (блокировании) денежных средств
или~иного имущества~\cite{rosfinmonitoring-lists,fz-115,fedsfm-mvk}.\footnote{Точные
  наименования перечней, формат публикации и~периодичность обновления сверяются
по~сайту Росфинмониторинга и~по~действующей редакции 115-ФЗ.}

Перечни Росфинмониторинга обновляются регулярно; после публикации обновлённого
перечня ОПДС обязан без~промедления, а~в~отдельных случаях не~позднее 24~часов
с~момента размещения информации, прогнать обновление и~применить меры к~клиентам,
попавшим в~перечень~\cite{rosfinmonitoring-lists,fz-115}.\footnote{Срок установлен
  115-ФЗ в~редакции Федерального закона от~15.12.2025~№~462-ФЗ; специальные
  правила исчисления сроков, ранее установленные приказами Росфинмониторинга,
отменены с~29.03.2026.}

\section{Международные перечни и их статус для ОПДС}\label{sec:sanctions-international}

Применение перечней ООН, ЕС, OFAC SDN и прочих определяется политикой ОПДС
по~риску, требованиями банков-корреспондентов; национальные перечни они не
подменяют~\cite{ofac-sanctions-list-service,eu-sanctions-info,unsc-consolidated-list}.

Транзакция в~валюте недружественной страны проходит через банк-корреспондент,
который применяет санкционные ограничения своей юрисдикции. Если ОПДС
формально не~обязан проверять контрагента по~OFAC SDN, отказ корреспондента
в~исполнении всё равно остановит операцию. ОПДС включает международные
перечни в~собственную проверку ради устойчивости канала расчётов, даже когда
национальный закон этого не~требует.

Executive Order~14114, подписанный 22~декабря 2023~года, дал OFAC право вводить
вторичные санкции против иностранного финансового института за~значительные
транзакции в~интересах российского военно-промышленного комплекса~\cite{ofac-eo-14114}.
15~января 2025~года OFAC впервые применил эту норму к~киргизскому Керемет-Банку
за~обслуживание подсанкционного Промсвязьбанка; банк потерял свыше~40\,\%
клиентских средств за~месяц после обозначения~\cite{ofac-keremet-bank-2025}.
Турецкие банки в~2024~году разорвали корреспондентские связи почти
со~всеми российскими организациями~\cite{kommersant-turkish-banks-2024}.
Многим иностранным банкам оказалось дешевле выйти из~российского направления
полностью, чем проверять каждую транзакцию на~риск вторичных санкций.

Помимо тройки OFAC SDN, EU Consolidated и~UN Consolidated, ОПДС
с~международными корреспондентами сверяется минимум с~четырьмя
самостоятельными перечнями.

\textbf{Council Regulation (EC) No 2580/2001.} Террористический список ЕС,
обновляемый Советом не~реже раза в~шесть месяцев~\cite{eu-2580-2001}. Узкий
перечень лиц, организаций и~групп, причастных к~терроризму, для~которых ЕС
применяет заморозку активов и~запрет на~предоставление финансовых услуг
отдельно от~общей санкционной массы.

\textbf{UK Sanctions List.} Единый британский список с~28~января 2026~года,
заменивший прежний OFSI Consolidated List of Financial Sanctions
Targets~\cite{uk-sanctions-list,uk-sanctions-single-list-2026}. Общий список объединяет финансовые,
иммиграционные, торговые и~транспортные санкции. Корреспонденты в~фунтах
проверяют контрагентов именно по~этому списку, а~не по~старому списку OFSI.

\textbf{DFAT Consolidated List.} Австралийский сводный перечень, который ведёт
Australian Sanctions Office при~Department of Foreign Affairs and Trade.
С~1~июля 2026~года подотчётные лица (\textit{reporting entities})
по~AML/CTF~Act~2006 обязаны устанавливать наличие клиента и~связанных лиц
в~списке как условие предоставления услуги~\cite{au-dfat-consolidated,austrac-tfs-reform}.

\textbf{MAS Targeted Financial Sanctions.} Сингапурский санкционный список
на~основе Financial Services and Markets Act~2022 и~Terrorism (Suppression of
Financing) Act~2002. Реализует резолюции Совета Безопасности ООН
и~собственные решения регулятора~MAS; финансовые институты Сингапура обязаны
замораживать активы лиц из~списка без~отдельного административного
акта~\cite{mas-tfs}.

Состав перечней зависит от~того, в~каких юрисдикциях ОПДС ведёт расчёты;
ограничения банка-корреспондента в~валюте перевода определяют его сильнее, чем
национальный закон самого ОПДС.

\subsection*{SDN и~SSI}

Минфин США ведёт семейство санкционных перечней с~разной правовой нагрузкой.

SDN (Specially Designated Nationals)~--- полная блокировка активов. Любая
операция через юрисдикцию США или~участников американской финансовой системы
с~SDN-лицом запрещена; американская организация обязана заморозить активы
и~отчитаться OFAC в~течение 10~рабочих дней. SSI (Sectoral Sanctions
Identifications)~--- секторальные санкции по~Executive Order~13662, охватывают
определённые секторы российской экономики через четыре
директивы~\cite{eo-13662,ofac-ssi-list}. Активы SSI-лица
не~блокируются; запрещены конкретные категории операций по~ограниченному
составу инструментов.

\textbf{Директива~1} (финансовый сектор)~--- новый долг сроком погашения
свыше~14~дней (для~долга, выпущенного с~28~ноября 2017~года; ранее~---
свыше~30~дней с~12~сентября 2014, свыше~90~дней с~16~июля 2014) или~новое
размещение акций (\textit{equity}).

\textbf{Директива~2} (энергетический сектор)~--- новый долг сроком
свыше~60~дней (с~28~ноября 2017; ранее~--- свыше~90~дней). Equity
не~ограничен.

\textbf{Директива~3} (оборонный сектор)~--- новый долг сроком свыше~30~дней.
Equity не~ограничен.

\textbf{Директива~4}~--- запрет на~поставку товаров, услуг (кроме финансовых)
и~технологий в~поддержку разведки или~добычи на~глубоководных, арктических
шельфовых и~сланцевых проектах в~России.

\highlight{Санкционный фильтр обязан различать классы}. Платёж \texttt{pacs.008}
в~евро с~SSI-контрагентом по~Директиве~1 проходит как обычная торговая операция;
тот же контрагент по~новому облигационному выпуску со~сроком погашения
31~день~--- нет. Сопоставление превращается из~бинарного в~матрицу
\enquote{лицо $\times$ тип операции $\times$ срок $\times$ продукт}. Нарушения
по~SSI состоят в~запрещённой механике расчёта при~формально допустимом контрагенте.

\subsection*{Правило~50\,\% OFAC}

Проверка по~имени проходит, но~сущность всё равно заблокирована~---
если её~50\,\% владельцев в~SDN-листе. Правило~50\,\% OFAC, в~действующей
редакции с~13~августа 2014~года, расширяет блокировку с~явных SDN-записей
на~любую сущность, у~которой одно или~несколько SDN-лиц владеют~50\,\%
или~более суммарно: прямо, косвенно через цепочку компаний-прокладок
или~агрегированно по~нескольким санкционным
владельцам~\cite{ofac-50-rule-guidance}.

Компания~X не~в~SDN-листе; её~30\,\% владеет SDN-лицо~A напрямую, ещё~25\,\%
владеет SDN-лицо~B через холдинг~Y (контролируемый~B на~80\,\%). Прямое
и~косвенное владение SDN-лицами суммируется: $30 + 25 \times 0{,}8 = 50\,\%$.
Компания~X заблокирована по~правилу~50\,\%, хотя в~SDN-листе её~нет, а~явных
владельцев~A и~B в~публичных реестрах может вообще не~быть.

OFAC не~публикует производный список таких компаний; обязанность построить граф
владения (\textit{ownership graph}) возложена на~ОПДС. Санкционная
проверка перестаёт сводиться к~сопоставлению имён (\textit{name matching})
и~требует той же работы по~построению графа, что и~KYB с~идентификацией
конечного бенефициарного владельца (\textit{UBO};
см.~раздел~\ref{sec:kyb-merchants}). Промышленные провайдеры (LexisNexis,
Refinitiv, Dow~Jones, Sayari) поддерживают индексы владения коммерчески;
в~собственной системе банка такой граф собирается на~базе данных KYB и~открытых
реестров бенефициарных владельцев. Правило~50\,\% BIS для~экспортного контроля
построено по~аналогичной схеме, но~применяется отдельно~--- эти правила
не~смешиваются.

\subsection*{14-й и~15-й пакеты ЕС}

В~2024~году санкции ЕС впервые ударили напрямую по~платёжной инфраструктуре,
а~не по~отдельным лицам и~компаниям. 14-й пакет ЕС от~24~июня 2024~года ввёл
запрет для~юрлиц ЕС вне России на~использование Системы передачи финансовых
сообщений Банка России (СПФС~--- российский аналог
SWIFT)~\cite{eu-14th-sanctions-2024}. Тот же пакет ввёл два новых механизма:
Annex~XLIV~--- список третьесторонних пользователей СПФС, подпадающих
под~санкции автоматически при~внесении; Annex~XLV~--- механизм листинга
финансовых институтов и~провайдеров криптоактивов из~третьих стран, вовлечённых
в~обход санкций.

Прежняя санкция накладывалась на~физлицо или~компанию~--- проверка шла
по~имени. Annex~XLIV и~XLV нацелены на~участие в~инфраструктуре~--- платёжный
шлюз, криптообмен, банк-корреспондент с~определённой ролью. Проверка учитывает
атрибуты самой транзакции~--- маршрут сообщения, банковский шлюз, использованный
криптомост,~--- а~идентификаторов сторон уже недостаточно. Для~банковских
транзакций структурированные поля приходят с~ISO~20022
(раздел~\ref{sec:sanctions-iso-20022}); для~криптомостов аналогичную роль играет
блокчейн-аналитика: сервисы Chainalysis, Elliptic, TRM~Labs
(раздел~\ref{sec:aml-blockchain-analytics}), разобранные для~финмониторинга,
размечают адреса и~трассируют переводы между блокчейнами (\textit{cross-chain})
в~реальном времени.

\section{Проверки по 860-П}\label{sec:sanctions-rf-modern}

\highlight{860-П встраивает санкционную проверку в~правила внутреннего контроля
(ПВК) кредитной организации как обязательную автоматизированную функцию}~\cite{cbr-860p}.

Кредитная организация обязана вести регулярную автоматизированную сверку
клиентов и~их контрагентов с~актуальной версией перечней Росфинмониторинга
и~решений МВК, фиксировать результат каждой сверки и~уровень совпадения
по~каждой паре. Без~автоматического движка с~журналом решений банк
не~укладывается ни в~24-часовой срок применения мер по~115-ФЗ, ни в~ежедневное
обновление санкционного индекса; ручной разбор такого объёма физически
невозможен. Полный разбор 860-П как ядра ПВК~--- в~главе~\ref{ch:aml}.

860-П заменяет 375-П, действовавшее с~2012~года; основное приращение~---
требование вести проверку как процесс с~зафиксированными уровнями совпадения.
Перечни и~решения МВК остаются исходными данными, но~к~ОПДС добавлены
требования к~самому процессу: автоматизированное сопоставление, регистрация
уровня совпадения, журнал снятия блока. Подходы вида \enquote{раз в~день
скачиваем CSV и~сверяем глазами} с~860-П несовместимы.

\section{Архитектура проверки}\label{sec:sanctions-screening-flow}

\highlight{%
  Сверка с~санкционными перечнями событийная: она идёт
  по~справочнику клиентов и~контрагентов и~запускается изменением записи
  или~обновлением перечня. На~горячем пути авторизации остаётся проверка двух
  флагов~--- у~клиента и~у~контрагента
}.

На~каждой записи справочника хранится результат последней сверки~--- флаг
состояния и~ссылка на~запись перечня при~совпадении. Авторизация читает эти
два поля (флаг клиента и~флаг контрагента) и~продолжает обычный путь, если
оба значения~--- \enquote{не~в~перечне}. Нечёткое сопоставление имён
в~горячем пути не~выполняется; вся вычислительная нагрузка вынесена в~события
сверки. Сверка запускается на~четырёх событиях.

\textbf{Загрузка справочника.} При~первом подключении к~актуальному перечню
каждая запись справочника прогоняется через нормализацию и~сравнение, флаг
записывается в~запись.

\textbf{Ввод нового клиента или~контрагента.} Запись прогоняется единожды
сразу после ввода; флаг сохраняется на~записи.

\textbf{Изменение существующей записи.} Перепрогоняется только эта запись;
флаг переписывается по~новому результату.

\textbf{Обновление перечня.} После~публикации обновлённого перечня
Росфинмониторинга весь справочник перепрогоняется против инкремента новой версии,
и~меры применяются к~попавшим в~перечень не~позднее 24~часов с~момента
размещения обновления~\cite{fz-115}. Под~860-П
(см.~раздел~\ref{sec:sanctions-rf-modern}) каждое такое событие сверки должно
дать запись в~журнале с~уровнем совпадения.

Записи с~нечётким совпадением выше порога автоматически попадают в~очередь
ручной эскалации; маршрут разбора и~критерии снятия блока описаны
в~разделе~\ref{sec:sanctions-false-positives}. Типовые ориентиры для~порога
Jaro-Winkler разобраны в~разделе~\ref{sec:sanctions-metrics}.

Ядро события сверки~--- нечёткое сопоставление пары строк. В~минимальной
реализации нормализация сворачивает имя к~канонической форме (удаление
  диакритики через одну из~форм нормализации Unicode~--- NFKD, Normalization
  Form Compatibility Decomposition~\cite{unicode-uax15}, нижний регистр,
отбрасывание пунктуации и~пробелов), после чего пара строк сравнивается
функцией отношения схожести строк. Превышение порога даёт совпадение; ниже~---
запись считается \enquote{не~в~перечне}. Рабочая система применяет более сложные
методы сопоставления:\nopagebreak%
\par\nobreak
\codefile{Python}{samples/ch27-sanctions/screening.py}

\subsection*{Метрики нечёткого сопоставления}\label{sec:sanctions-metrics}

Сопоставление имён в~рабочей системе сводится к~выбору метрики и~порога. Почти
все сценарии санкционной проверки покрываются четырьмя метриками:
редакторская дистанция Левенштейна, расширение Дамерау~--- Левенштейна
с~транспозицией, мера Jaro и~Jaro-Winkler с~бонусом за~общий префикс.

\textbf{Levenshtein.} Минимальное число вставок, удалений и~замен,
превращающих одну строку в~другую~\cite{levenshtein-1965}. Чувствительна
к~длине строки и~к~любому изменению независимо от~его природы.

\textbf{Damerau-Levenshtein.} Расширение Левенштейна с~операцией транспозиции
соседних символов: перестановка двух соседних букв стоит~1
вместо~2~\cite{damerau-1964}. Отделяет опечатку клавиатуры от~двух
независимых правок.

\textbf{Jaro.} Доля совпадающих символов в~окне ширины
$\lfloor\max(|a|,|b|)/2\rfloor-1$ со~штрафом за~транспозиции. Значение
в~$[0, 1]$, в~отличие от~целочисленной редакторской
дистанции~\cite{jaro-1989}.

\textbf{Jaro-Winkler.} Модификация Jaro: бонус за~общий префикс длиной
до~четырёх символов с~коэффициентом $p=0{,}1$~\cite{winkler-1990}. Имена
с~совпадающим началом получают усиление; ошибки чаще встречаются в~конце
имени, реже в~начале.

В~Python все четыре метрики, а~также фонетические коды Soundex и~Metaphone,
берутся готовыми из~пакета \texttt{jellyfish}~\cite{jellyfish}.
На~типичных кириллических парах метрики расходятся так
(табл.~\ref{tab:sanctions-metrics}).

\begin{table}[H]
  \begingroup
  \footnotesize
  \setlength{\tabcolsep}{6pt}
  \renewcommand{\arraystretch}{1.2}
  \begin{tabularx}{\textwidth}{
      @{}B{0.34\textwidth} Y Y Y Y@{}
    }
    \toprule
    \headrow
    \thd{Пара} & \thd{Lev} & \thd{DL} & \thd{Jaro} & \thd{JW} \\
    \midrule
    IVANOV~--- IVANOV       & 0 & 0 & 1,000 & 1,000 \\
    IVANOV~--- IVANOFF      & 2 & 2 & 0,849 & 0,910 \\
    IVANOV~--- YVANOV       & 1 & 1 & 0,889 & 0,889 \\
    IVANOV~--- IVAN-OFF     & 3 & 3 & 0,819 & 0,892 \\
    IVANOV~--- IVAONV       & 2 & 1 & 0,944 & 0,961 \\
    \bottomrule
  \end{tabularx}
  \endgroup
  \caption{Метрики на~парах с~разной природой
  расхождения.}\label{tab:sanctions-metrics}
\end{table}

Только Damerau-Levenshtein отделяет транспозицию NO~$\leftrightarrow$~ON (пара
IVAONV~--- IVANOV) от~двух независимых правок и~считает её одной операцией,
а~не~двумя.
Jaro-Winkler ценит общий префикс \textbf{IVAN}: пара IVANOV~--- IVANOFF
получает прирост $+0{,}061$ против чистого Jaro благодаря бонусу за~четыре
совпадающих начальных символа. Пара IVANOV~--- YVANOV того же бонуса
не~получает, потому что расходится на~первом символе.

Отраслевой порог Jaro-Winkler в~санкционной проверке~--- около $0{,}85$.
У~Levenshtein и~Damerau-Levenshtein порог задаётся относительно длины имени:
одна-две правки на~короткое имя.

Фонетические алгоритмы Soundex (Russell 1918~\cite{russell-soundex})
и~Metaphone (Philips 1990~\cite{philips-metaphone}) кодируют имя в~короткую
строку, и~сопоставление сводится к~сравнению кодов. Оба алгоритма построены
под~английскую фонологию и~для~кириллицы без~предварительной транслитерации
выдают мусор. Применяются поверх транслитерированных форм; в~санкционной
проверке работают как дополнение к~метрикам редакторской дистанции, не как
основной фильтр.

\subsection*{Стандартизация транслитерации}\label{sec:sanctions-translit}

\highlight{Стандарт транслитерации, который применяет CMS, определяет, что
попадает в~очередь ручного разбора}. Один и~тот же человек в~разных формах
транслитерации даёт разные строки, и~Jaro-Winkler ставит им разные значения~---
иногда выше порога, иногда ниже.

Для российских имён в~ходу три стандарта.

\textbf{ISO 9:1995} (тождествен ГОСТ 7.79-2000 системе~А)~--- одна латинская
буква с~диакритикой за~каждую кириллическую (Ë, Ž, Č, Š, Ŝ, Û, Â). Самые
короткие строки, никаких контекстных правил; полностью обратимо без~потери
буквы~\cite{iso-9}.

\textbf{BGN/PCGN 1947}~--- стандарт США (Board on Geographic Names)
и~Великобритании (Permanent Committee on Geographical Names) для~документов
и~карт. Шипящие раскрываются в~диграфы (KH, TS, CH, SH, SHCH); сохраняется
диакритика для~ё. Единственное контекстное правило: Е и~Ё транслитерируются
как ye, yë в~начале слова, после гласных и~после й, ъ, ь, иначе e,
ë~\cite{bgn-pcgn}. Обратимость частичная.

\textbf{ICAO Doc 9303 part 3}~--- стандарт ИКАО для~машиночитаемой зоны (MRZ)
паспорта. Только ASCII A-Z, без~диакритики и~без~контекстных правил; ь
не~передаётся, ё транслитерируется как e (сливается с~е)~\cite{icao-9303}.
Обратимости нет.

На имени ЮРИЙ расхождение становится критическим: между YURIY (BGN/PCGN)
и~IURII (ICAO) Jaro выдаёт $0{,}733$~--- ниже типового порога $0{,}85$. Человек
с~внутренним документом, оформленным по~BGN/PCGN, при~проверке не~сопоставляется
с~санкционным перечнем, где это же имя стоит в~ICAO-форме, без~специальной
нормализации обеих сторон к~единому стандарту.

Это расхождение видно в~карточной эмиссии: один банк требует от~клиента вписать
имя \enquote{как в~заграничном паспорте}~--- то есть ровно по~ICAO Doc 9303,
и~ответственность за~разночтение в~международной идентификации переносится
на~сам документ. Другой банк генерирует имя сам по~внутренней таблице, иногда
не~совпадающей ни с~одним международным стандартом. ЕВГЕНИЙ Иванов получает в~разных банках
разные формы для~своей карты: в~одном банке EVGENIY по~BGN/PCGN,
у~соседа~--- EVGENII по~ICAO, у~третьего~--- EVGENIJ по~ISO 9. Каждая форма
правильна по~своему стандарту; санкционная проверка зависит от~того, какую
из~трёх видит CMS банка и~какие формы внесены в~санкционный индекс. Различие
в~одном символе на~коротком имени даёт перепад Jaro-Winkler в~$0{,}05$--$0{,}08$,
ровно на~типовом пороге $0{,}85$.

Все три стандарта реализованы в~Python-пакете \texttt{iuliia}~\cite{iuliia}.

\begin{table}[H]
  \begingroup
  \footnotesize
  \setlength{\tabcolsep}{6pt}
  \renewcommand{\arraystretch}{1.2}
  \begin{tabularx}{\textwidth}{
      @{}B{0.20\textwidth} Y Y Y@{}
    }
    \toprule
    \headrow
    \thd{Исходное} & \thd{ISO 9:1995} & \thd{BGN/PCGN} & \thd{ICAO 9303} \\
    \midrule
    ЕВГЕНИЙ    & EVGENIJ    & YEVGENIY    & EVGENII \\
    ЕКАТЕРИНА  & EKATERINA  & YEKATERINA  & EKATERINA \\
    ЮРИЙ       & ÛRIJ       & YURIY       & IURII \\
    ЦАРЁВ      & CARËV      & TSARËV      & TSAREV \\
    ЩУКИН      & ŜUKIN      & SHCHUKIN    & SHCHUKIN \\
    ФЁДОРОВ    & FËDOROV    & FËDOROV     & FEDOROV \\
    \bottomrule
  \end{tabularx}
  \endgroup
  \caption{Транслитерация русских имён по~трём
  стандартам.}\label{tab:sanctions-translit}
\end{table}

Тот же эффект виден на~китайских, арабских и~корейских именах. Одну и~ту же
китайскую фамилию пиньинь записывает как \textit{Zhang}, Уэйда-Джайлз~---
\textit{Chang}, кантонская романизация~--- \textit{Cheung}; одно арабское имя
становится \textit{Mohammad}, \textit{Mohamed}, \textit{Muhammad} или~\textit{Mohammed}
в~зависимости от~стандарта латинизации (UNGEGN, ALA-LC, BGN/PCGN). OFAC SDN ведёт
записи в~собственных формах, коммерческие санкционные индексы партнёров~--- в~своих.
Совпадение формы имени между CMS и~санкционным индексом не~гарантировано ни
для~одной из~этих систем письма.

\section{{ISO}~20022 и~санкционная проверка}\label{sec:sanctions-iso-20022}

\highlight{22~ноября 2025~года SWIFT завершил период сосуществования форматов:
  все трансграничные платежи между финансовыми институтами обязаны использовать
ISO~20022~MX}~\cite{iso-20022-cbpr-nov-2025}. MT103 (клиентский перевод), MT202
(межбанковский перевод) и~MT940 (выписка) заменены на~\texttt{pacs.008},
\texttt{pacs.009} и~\texttt{camt.053} соответственно.

MT-сообщения хранили данные сторон в~неструктурированном текстовом блоке~---
одна строка \texttt{50K:} с~именем и~адресом, разделёнными переносом. Имя
контрагента, фамилия, юридический адрес и~страна оказывались в~одном поле
без~формальной разметки. Санкционный фильтр извлекал имя эвристиками и~работал
на~\enquote{сырых} строках. MX-сообщения (CBPR+~--- Cross-Border Payments and
Reporting Plus) задают структурированные поля~\cite{iso-20022-cbpr-plus}:

\begin{enumerate}
  \item Nm~--- имя;
  \item PstlAdr~--- почтовый адрес с~townName, country code, subdivision;
  \item LEI~--- Legal Entity Identifier для~юрлиц;
  \item Id/OrgId/AnyBIC~--- идентификатор организации.
\end{enumerate}

Проверка работает с~отдельными атрибутами, а~не с~единой строкой, как раньше.
Точность сопоставления выросла: имя и~адрес разделены, код страны в~отдельном
поле проверяется по~стране-юрисдикции, LEI даёт детерминированный ключ
сопоставления для~юридических лиц с~зарегистрированной идентификацией.
Одновременно вырос объём ответственности ОПДС~--- неполные
или~некорректные структурные поля больше не~прикрываются \enquote{общим
текстом}; банк-отправитель обязан проверить наличие обязательных атрибутов
до~передачи сообщения дальше.

Переходный период гибридных адресов завершается 14~ноября 2026~года
(SWIFT~SR~2026). До~этой даты допустим \enquote{гибридный} формат: townName
и~country code обязательны как~структурированные поля, остальные элементы
адреса допустимы в~неструктурированных строках (\texttt{AdrLine}, до~двух
строк по~70~символов). После 14~ноября 2026~года неструктурированные
почтовые адреса в~платёжных сообщениях CBPR+ отбиваются сетью (NAK);
исключения сохранены для~выписок и~уведомлений
(\texttt{camt.052/053/054}, \texttt{admi.024}, \texttt{camt.025},
\texttt{camt.060}).

Около 300 российских банков формально сохраняют SWIFT-подключение,
и~трансграничные платежи для~несанкционированных корпоративных клиентов идут
через них по~ограниченному набору официально документированных маршрутов. Российские
дочки западных банков~--- АО~Райффайзенбанк (Raiffeisen Bank International)
и~АО~ЮниКредит Банк (UniCredit)~--- под~санкции ЕС не~подпадают и~остаются
основным каналом евро-расчётов; ЦБ~РФ держит их~в~списке системно
значимых. Газпромбанк до~21~ноября 2024~года был основным каналом
энергетических платежей под~исключением ЕС для~нефти и~газа; OFAC включил
его~в~список SDN 21.11.2024 по~Executive~Order~14024, и~после истечения
\textit{wind-down} по~General~Licence~8L 12~марта 2025~года долларовый коридор
для~него закрылся. Регламенты ЕС евро-платежи Газпромбанку не~запрещают, но~круг
готовых корреспондентов после~OFAC-обозначения резко сузился. Поверх
адресных банковских санкций OFAC и~ЕС держат набор узких типизированных
исключений: General~Licence~6D OFAC разрешает поставки в~РФ
сельскохозяйственной продукции и~медицинских товаров; GL~55C~--- расчёты
по~Sakhalin-2 LNG до~28.06.2025; GL~53A~--- дипломатические операции
с~Газпромбанком; Регламент ЕС~833/2014 в~преамбуле §44 исключает поставки
Росатома для~АЭС~Paks~II в~Венгрии. Сам Росатом как организация под~целевой
запрет ЕС не~подпадает~--- индивидуальные дочерние компании указываются
в списках точечно.

\section{Wire stripping}\label{sec:sanctions-wire-stripping}

\highlight{\textit{Wire stripping}~--- удаление или~подмена корреспондентским
  банком полей \texttt{50K} (отправитель), \texttt{59} (бенефициар),
  \texttt{56}/\texttt{57} (посредники) в~исходящем SWIFT MT-сообщении, в~которых
  видно имя или~BIC санкционного контрагента, чтобы следующий корреспондент
на~маршруте через США или~ЕС не~получил совпадения на~OFAC- или~EU-фильтре}.
Делает это сам корреспондент, через которого санкционное лицо отправляет
валютный платёж: переписывает поля, заменяет имя нейтральным
юрлицом из~той же группы и~направляет \enquote{очищенное} сообщение дальше
по~маршруту. Корреспондент получает комиссию с~санкционного клиента,
а~проверка входящей операции видит сообщение, неотличимое от~законного
перевода.

Standard~Chartered Bank с~2001 по~2007~год систематически удалял санкционные
идентификаторы (SWIFT/BIC иранских контрагентов) из~MT103 и~MT202
на~корреспондентском маршруте через Нью-Йорк, чтобы пройти OFAC-проверку
американских клиринговых банков. NYDFS в~2012~году выписал штраф \$340\,млн
и~опубликовал детальный анализ схемы~\cite{nydfs-stanchart-wire-stripping-2012}.
Урегулирование 2012~года практику не~прекратило: в~2019~году тот же банк
по~совокупному соглашению с~NYDFS, DOJ, OFAC, ФРБ Нью-Йорка и~британским FCA
выплатил ещё~\$1,1\,млрд за~нарушения по~Бирме, Кубе, Ирану, Судану и~Сирии
в~период 2008--2014~годов через дубайский филиал и~онлайн-платформу
\enquote{Straight-to-Bank}~\cite{nydfs-stanchart-2019}.

BNP~Paribas повторил тот же сценарий в~больших масштабах. С~2004 по~2012~год
банк использовал wire stripping и~непрозрачную маршрутизацию через сеть
корреспондентов для~обхода санкций США против Судана, Ирана и~Кубы.
Deutsche~Bank в~2015~году получил аналогичные предписания от~NYDFS
на~\$258\,млн.

Структурная валидация схемы сама по~себе wire stripping не~исключает:
банк-корреспондент формирует исходящее MX-сообщение по~собственному входу
(инструкция приходит к~нему по~любому транспорту), и~парсер получающего банка
увидит структурно валидное \texttt{pacs.008}\slash\texttt{pacs.009} независимо
от~того, реальный отправитель в~поле \texttt{Nm} или~оболочка. Закрывают вектор
атаки три отдельных правила поверх транспортной схемы.

\textbf{Cover payment transparency.} MT202COV (введён SWIFT 21~ноября 2009~года)
и~ISO~20022 \texttt{pacs.009} со~связанной \texttt{pacs.008} требуют, чтобы
межбанковское cover-сообщение нёсло данные \textit{исходного клиентского перевода}.
До~введения MT202COV корреспондент мог провести MT103 с~санкционным клиентом
через свои внутренние счета или~сателлитный банк, а~в~Нью-Йорк
для~долларового клиринга отправить отдельный MT202 без~клиентских данных~---
американский клиринг видел только межбанковскую часть маршрута и~не~получал
ничего для~OFAC-фильтра. На~этом разрыве BNP~Paribas строил суданские
платежи 2002--2007 через сеть сателлитных банков в~Африке и~на~Ближнем
Востоке.

\textbf{Travel Rule.} Рекомендация~16 FATF и~для~ЕС Регламент~(EU)
2023/1113 обязывают каждое платёжное сообщение нести полные данные отправителя
и~бенефициара (имя, счёт, адрес или~идентификатор); банк-получатель должен
отбивать сообщения с~неполной информацией~\cite{fatf-recs,eba-travel-rule-2024}.
Скрытие имени превращается в~отдельное регулятивное нарушение, а~не~в~незаметное
удаление строки.

\textbf{Машиночитаемый аудит.} Код страны, LEI и~BIC в~отдельных полях позволяют
надзорному органу или~банку-корреспонденту прогнать массив сообщений ex~post
по~детерминированному фильтру (страна\,=\,IR, BIC из~SDN-списка) без~разбора
свободного текста. Поиск становится дешёвым, и~систематическое применение схемы
ловится по~статистике потоков, а~не~по~отдельному сообщению.

Эти правила закрывают именно \enquote{очистку} существующих полей. Подмена
имени реального клиента на~юрлицо из~той~же группы остаётся
отдельным вектором: банк-эмитент заполняет обязательные поля синтаксически
правильным, но~ложным значением. Схема не~отличает подлинного отправителя
от~подложного, и~против такого сценария работает уже~KYC
на~стороне банка-эмитента (см.~гл.~\ref{ch:kyc-kyb}).

LEI обходит проблему транслитерации, но не~решает её: совпадение
проверяется посимвольно по~20-значному коду, имя в~процедуре
не~участвует. Само поле \enquote{legal name} в~LEI-реестре хранится
в~исходной графике (для~российских юрлиц~--- кириллицей),
а~ASCII-транслитерация опциональна и~ставится локальным регистратором
по~усмотрению, так что многовариантность имени сохраняется и~внутри
LEI-данных. Условие работы~--- LEI присутствует и~у~отправителя
в~MX-сообщении, и~у~санкционного субъекта в~перечне; OFAC~SDN
и~EU~Consolidated проставляют LEI не~для~всех записей, а~физлица
в~LEI не~входят по~определению.

\section{Ложные срабатывания и маршрут разбора}\label{sec:sanctions-false-positives}

Ложноположительные совпадения (гл.~\ref{ch:fraud}) попадают в~отдельную
очередь: система помечает запись справочника как требующую проверки, аналитик
смотрит контекст и~либо снимает блок, либо эскалирует случай. Маршрут должен
быть формализован, иначе одинаковые срабатывания приводят к~разным решениям.
Распространённые имена и~совпадения между кириллицей и~латиницей дают много
шума, и~значительная доля работы аналитика~--- быстрый разбор заведомо ложных
срабатываний~\cite{ofac-sanctions-faq}. Логика сопоставления и~пороги
срабатывания настраиваются исходя из~собственной оценки риска ОПДС,
зафиксированной в~правилах внутреннего контроля по~115-ФЗ.

Аналитик получает карточку совпадения с~атрибутами операции (страна, сумма,
продукт, история клиента), оценкой схожести по~основной метрике и~ссылкой
на~запись в~перечне. Регламент разрешает аналитику самому снять блок
с~обоснованием в~журнале и~отпустить транзакцию; эскалация в~службу комплаенс
обязательна, если совпадение касается санкционных перечней с~правовой нагрузкой
(статья~6 115-ФЗ, OFAC SDN, EU Consolidated), если клиент уже фигурировал
в~эскалациях или~если запрос инициирован банком-корреспондентом.

Журнал решений~--- такая же часть санкционного контроля, как и~сам алгоритм
сопоставления. Регулятор и~аудит проверяют у~ОПДС воспроизводимую процедуру
принятия решения и~след каждого снятия блока; срабатывание метрики
на~конкретной паре само по~себе ничего не~доказывает. Без~журнала
отсутствие нарушений нельзя доказать.
